\section{Introduction}
\label{sec:intro}

Tensor network models have long represented the state of the art in modeling complex quantum systems~\citep{white1992, fannes1992, orus2019}, but have only recently been utilized as models for machine learning~\citep{novikov2015, cohen2016, stoudenmire2016, novikov2017, han2018, stoudenmire2018, cheng2019}. In contrast to neural networks, tensor networks forgo the use of nonlinear activation functions, relying instead on multiplicative interactions to capture complex correlations within data. This gives tensor networks a convenient mathematical structure suitable for proving general theoretical results, such as the separation in expressivity between almost all deep tensor networks and their shallow counterparts~\citep{cohen2016}. However, these distinctive mathematical properties have yet to be leveraged for the development of new \emph{operational} abilities, which would 
give more practical reasons for the wider adoption of tensor network models in real-world machine learning tasks.
 
In this work we apply a recurrent tensor network, the \emph{uniform matrix product state}~(u-MPS), to the task of probabilistic sequence modeling, and identify several novel abilities of u-MPS regarding their evaluation and generative capabilities. Despite its recurrent nature, we show that sequential inputs to u-MPS can be processed in a highly parallel manner, with sequences of length $n$ being evaluated in parallel time $\bigO(\log n)$. While the difficulty of parallelizing deep recurrent neural networks (RNNs) has previously motivated the development of non-recurrent architectures for sequence processing tasks (e.g. \citep{gehring2017, vaswani2017}), our finding shows that recurrent tensor networks represent another means of achieving greater parallelism.

We further show that u-MPS models are endowed with surprising generative capabilities closely tied to the structure of regular expressions (regex). While standard autoregressive models are constrained to generate sequences in a stream-like fashion conditioned on some prompt, we find that u-MPS can sample from a wide variety of distributions defined by conditioning regular expressions $R$. Our sampling algorithm efficiently produces unbiased samples from the probability distribution learned by the u-MPS, conditioned on the output sequence matching a given regular expression $R$. Standard autoregressive sampling follows from the choice $R = p\Sigma^*$ (for $p$ a prefix string and $\Sigma^*$ the regex matching all sequences), but other special cases include fill-in-the-blank sampling ($R = p\Sigma^*s$, for suffix $s$), as well as the generation of samples constrained to contain some target phrase $t$ ($R=\Sigma^*t\Sigma^*$).

Besides permitting the generation of sequences with rich internal structure, these techniques enable novel forms of regularization, where a u-MPS model can be penalized or incentivized during training to generate strings matching some target pattern. Such regularization can be applied in a variety of challenging tasks, including automatic code generation and mitigating gender bias in language models. Experiments on synthetic and real structured text datasets confirm these novel parallelism, sampling, and regularization benefits, and show u-MPS able to successfully generalize non-local correlations present in small strings to sequences of significantly greater length.

\paragraph{Summary of Contributions}
We give the first implementation of a u-MPS for probabilistic sequence modeling, and uncover several remarkable capabilities of this model\footnote{The code used to produce our results can be found at \url{https://github.com/jemisjoky/umps_code}.}. The absence of nonlinear activation functions in the u-MPS allows us to utilize a parallel evaluation method during training and inference. We develop new techniques linking the structure of u-MPS ``transfer operators'' to that of regular expressions, which in turn enables a flexible recursive sampling algorithm and novel forms of regularization for the u-MPS. We expect these techniques to open up significant new research directions in the design of sequential generative models, with language modeling being a particularly promising domain.

\paragraph{Related Work}
Notable previous applications of tensor networks in machine learning include compressing large neural network weights~\citep{novikov2015}, proving separations in the expressivity of deep vs shallow networks~\citep{cohen2016}, and for supervised~\citep{stoudenmire2016, novikov2017, glasser2018} and unsupervised~\citep{han2018, stoudenmire2018, cheng2019} learning tasks. Of particular relevance is~\citep{stokes2019}, where (non-uniform) MPS were trained as generative models for fixed-length binary sequences using the density matrix renormalization group (DMRG) algorithm. A diverse range of tensor network architectures have also been proposed as theoretical tools for modeling and understanding natural language, such as~\citep{pestun2017, coecke2010, gallego2017, degiuli2019}. The completely positive maps employed in our sampling algorithm are similar to those used in hidden quantum Markov models (HQMM)~\citep{monras2010,srinivasan2018}, and can likewise be interpreted using concepts from quantum information theory.

This work can be seen as a continuation of~\citep{pestun2017a}, where u-MPS were introduced from a theoretical perspective as a language model, but without the parallelization, sampling, or experimental results given here. Our sampling algorithm is a significant generalization of the fixed-length Born machine algorithm introduced in~\citep{han2018}~(which in turn follows that of~\citep{ferris2012}), and by virtue of the recurrent nature of u-MPS, permits the generation of discrete sequences of arbitrary length. The u-MPS model is equivalent to quadratic weighted finite automata~\citep{bailly2011} and (to a lesser extent) the norm-observable operator model (NOOM)~\citep{zhao2010}, and is also an example of a linear (second-order) RNN~\citep{rabusseau2019}. Benefits of linear (first-order) RNNs for parallelization and interpretability were described in~\citep{martin2018, foerster2017input}.

A key difference from previous works is the general techniques developed for evaluating and sampling from regex-structured probability distributions, which to the best of our knowledge are completely new. These techniques apply not only to u-MPS, but also to a broad family of models with similar internal structure, such as weighted finite automata (WFA)~\citep{droste2009}, hidden Markov models (HMM)~\citep{rabiner1986}, and predictive state representations~\citep{littman2002}. We consequently expect the algorithms developed here for regex sampling, regularization, and parallel evaluation to immediately generalize to any of these models which parameterize valid probability distributions.

